\documentclass[handout]{beamer} %Remove handout to reenable pauses
\usetheme{Berlin}
\usepackage{array}
\usepackage{color}
\usepackage{graphicx}
\usepackage{tikz}
\usepackage{tikz-cd}
\usepackage{wrapfig}
\usepackage{amsmath}

\usetikzlibrary{matrix, shapes, trees, arrows, decorations.pathmorphing, fit}
\usepackage[utf8]{inputenc}

\DeclareMathOperator{\K}{K}
\DeclareMathOperator{\G}{G}
\DeclareMathOperator{\SL}{SL}
\DeclareMathOperator{\Sp}{Sp}
\DeclareMathOperator{\GL}{GL}
\DeclareMathOperator{\E}{E}
\DeclareMathOperator{\St}{St}
\DeclareMathOperator{\Cent}{Cent}
\DeclareMathOperator{\sr}{sr}
\DeclareMathOperator{\rk}{rk}
\DeclareMathOperator{\Img}{Im}
\DeclareMathOperator{\Umd}{Umd}
\DeclareMathOperator{\Ker}{Ker}
\DeclareMathOperator{\coker}{coker}
\newcommand{\catname}[1]{{\normalfont\textbf{#1}}}
\newcommand{\Groups}{\catname{Groups}}
\newcommand{\rA}{\mathsf{A}}
\newcommand{\rB}{\mathsf{B}}
\newcommand{\rC}{\mathsf{C}}
\newcommand{\rD}{\mathsf{D}}
\newcommand{\rE}{\mathsf{E}}
\newcommand{\rF}{\mathsf{F}}
\newcommand{\rG}{\mathsf{G}}
\newcommand{\StH}{\widehat{H}}
\DeclareMathOperator{\UU}{U}

\theoremstyle{definition}
\newtheorem*{thm}{Theorem}
\newtheorem*{prop}{Proposition}
\newtheorem*{lem}{Lemma}
\newtheorem*{exm}{Example}
\newtheorem*{dfn}{Definition}
\newtheorem*{que}{Question}
\newtheorem*{conj}{Conjecture}
\newtheorem*{ans}{Answer}
\newtheorem*{rem}{Remark}
\newtheorem*{cor}{Corollary}
\newtheorem*{conc}{Conclusion}
\newtheorem*{fac}{Fact}
\newtheorem*{ide}{Idea}
\newtheorem*{prb}{Problem}

\def\hnode{*[o]{\circ}}
\def\lnode#1{*[o]{\composite{!{\circ}*h!<0pt,-1em>{\scriptstyle #1}}}}
\def\lnoded#1{*[o]{\composite{!{\circ}*h!<0pt,0.5em>{\scriptstyle #1}}}}
\def\lbnode#1{*[o]{\composite{!{\bullet}*h!<0pt,-1em>{\scriptstyle #1}}}}
\def\harr#1{\ar@{-}[r]_{\displaystyle #1}}
\def\hard{\ar@{-}[d]}
\def\harl#1{\ar@{-}[l]^{\displaystyle #1}}

\newcommand\restr[2]{{\left.\kern-\nulldelimiterspace #1\vphantom{\big|}\right|_{#2}}}

\tikzset{
  root/.style={
    circle,
    draw,
    fill=black,
    inner sep=1.5pt
  },
  highlighted/.style={
    fill=green!60!black
  },
  zeroroot/.style={
    fill=white
  },
  levi/.style={
    draw,
    rounded corners,
    inner sep=3pt
  },
  dottededge/.style={
    dotted,
    thick
  }
}

\title{Suslin decompositions of Steinberg groups and the $K_2$-analogue of Bass--Quillen conjecture}
\author{Sergei Sinchuk}
\date{6 March 2026}

\AtBeginSection[]{}

%%%%%%%%%%%%%%%%%%%%%
\begin{document}

\begin{frame}
\titlepage
\end{frame}

\section{BQ and $\mathbb{A}^1$-homotopy}

\begin{frame} \tableofcontents[currentsection] \end{frame}

\begin{frame}{Bass--Quillen Conjecture}

\textbf{Conjecture (Bass--Quillen).}

Let $R$ be a \emph{regular Noetherian ring}. Then every finitely generated projective module over the polynomial ring
\[R[t_1,\dots,t_n]\]
is \emph{extended from $R$}.
\pause

Equivalently: if $P$ is a finitely generated projective $R[t_1,\dots,t_n]$-module, then
\[P \cong P_0 \otimes_R R[t_1,\dots,t_n]\]
for some finitely generated projective $R$-module $P_0$.
\pause

\medskip

\textbf{Geometric formulation.}

For $X=\mathrm{Spec}(R)$ regular, every vector bundle on $X \times \mathbb{A}^n$ is pulled back from $X$, so $\mathrm{Vect}_n(X) = \mathrm{Vect}_n(X \times \mathbb{A}^n)$
\end{frame}

\begin{frame}{Known Cases}

The conjecture is known in the following situations:

\begin{itemize}
  \item $R$ a Noetherian regular comm. ring containing a field (Lindel--Popescu theorem 1986);
  \pause
  \item $R$ regular Noetherian of dimension $\leq 2$ (Quillen--Suslin--Murthy);
  \pause
  \item $R$ smooth over $k$, $k$ Noetherian regular over a Dedekind domain $D$ with perfect residue fields (Asok--Hoyois--Wendt 2015).
\end{itemize}
\pause
Key tools include Popescu's \emph{general N\'eron desingularization}.
\pause

Regular morph. = flat morph. + regular geom. fibers;
\pause

Smooth morph. = Regular morph. + finite type.

\end{frame}

\begin{frame}{$\mathbb{A}^1$-homotopy theory (Morel--Voevodsky, late 90s)}
 \[\mathcal{H}_{k}^{\mathbb{A}^1} = \mathrm{Ho}\left( (\mathcal{S}_\mathrm{Nis} \cup \mathcal{S}_{\mathbb{A}^1})^{-1} \mathrm{Fun}(\Delta^{op}, \mathrm{Presh}(\catname{Sm}_k)) \right)\]
 \pause
 \[\mathcal{S}_{\mathbb{A}^1} = \{ X \times_k \mathbb{A}^1_k \to X\}, \]
 \pause
 \[\mathcal{S}_\mathrm{Nis}\text{ -- covering sieves in Nisnevich topology} \]
 \pause
 Nisnevich topology is generated by the following squares
 \[ \begin{tikzcd}[ampersand replacement=\&, column sep=1.1em] \& Y \arrow{d}{p\text{ -{\'e}t}} \& (Y \setminus p^{-1}(U))_\mathrm{red} \arrow{d}{\cong} \\ U \arrow[r, hookrightarrow] \& X \& (X \setminus U)_\mathrm{red} \end{tikzcd}\]
 \pause

 \begin{thm}[Asok--Hoyois--Wendt, 2015]
  Let $X = \mathrm{Spec}(R)$ one has $\mathcal{H}_{k}^{\mathbb{A}^1}(X, \mathrm{Gr}_{r, \infty}) = \mathrm{Vect}_{r}(X) \equiv \mathrm{K}_{0, r}(R).$
 \end{thm}

\end{frame}

\begin{frame}{Topological picture}

$X$ paracompact Hausdorff top. space.
\pause

$\mathrm{Vect}^{top}_{r}(X)$ - set of top. real vect. bundles over $X$.
\pause

It obviously satisfies $\mathrm{Vect}^{top}_{r}(X) \cong \mathrm{Vect}^{top}_{r}(X \times \mathbb{R})$.
\pause

\begin{thm}[Steenrod 1951]
 \(\mathrm{Vect}^{top}_{r}(X) = \mathcal{H}(X, \mathrm{Gr}^{top}_{r}(\mathbb{R}^\infty)) = \mathcal{H}(X, BO_r) = \mathcal{H}(X, BGL_r).\)
\end{thm}
\pause

Here $O_r$ and $GL_r$ are understood as real Lie groups.
\pause

Replacing $GL_r$ with $Sp_{2r}$ (resp. anisotropic $U(r)$) classifies quaternionic rank $r$ (complex rank $r$) top. vect. bundles.
\end{frame}

\begin{frame}{Classification of torsors}

Let $G$ be an algebraic $k$-group. The bar construction $BG_\bullet$ is the Yoneda-image of the simplicial scheme given by bar-construction.
\pause

\begin{thm}[Asok--Hoyois-Wendt, 2015] If $k$ is a field and $G$ is an isotropic reductive $k$-group, then for every smooth affine $k$-scheme $X$ there is a functorial bijection
\[
H^1_{\mathrm{Nis}}(X,G) \;\cong\; \mathcal{H}(X,BG)_{\mathbb{A}^1}.
\]
\end{thm}
\pause
$H^1_{\mathrm{Nis}}(X,G)$ is the set of $G$-torsors on $X$, loc. triv. in Nisnevich topology.
\end{frame}

\begin{frame}{$H^1_{\mathrm{Nis}}(X,G)$ }

\textbf{Definition}
A \emph{$G$-torsor} over $X$ is an $X$-scheme $P \to X$ together with a (right) action
\( P \times_X G \to P \) such that

\begin{itemize}
\item the action is free and transitive on fibers, equivalently
\[ P \times_X G \;\xrightarrow{\;\sim\;}\; P \times_X P, \quad (p,g) \mapsto (p,pg), \]
\pause
\item $P$ is locally trivial, i.e. there exists a Nisnevich cover $\{U_i \to X\}$
such that \[ P \times_X U_i \cong G \times_X U_i . \]
\end{itemize}
\pause

\textbf{Quadratic case:} ($2 \in \mathcal{O}_X^\times$) For $G = O_r$, $H^1_{\mathrm{Nis}}(X,O_n)$ classifies quadratic spaces $(E,q)$ of rank $r$ s.t. for every point $x \in X$ after passing to $\mathcal{O}_{X,x}^h$, the quadratic form becomes split.
\end{frame}

\begin{frame}{Analogue of Bass-Quillen for $\K_1$}
 Compute $\mathcal{H}_{k,*}^{\mathbb{A}^1}(X_+, G) = \mathcal{H}_{k}^{\mathbb{A}^1}(X, G)$ (just drop $B$)?
\pause

 \begin{thm}[Stavrova, 2011, 2019]
 Let $G$ be a reductive group scheme over $k$ of isotropic rank $\geq 2$.
 Let $A$ be a regular ring containing a field $k$. Then there is a natural isomorphism
 $K_1^G(A) \cong K_1^G(A[X]).$
 Hence $K_1^G(A) = \mathcal{H}_{k,*}^{\mathbb{A}^1}(X_+, G)$.
 \end{thm}
\pause

 \begin{thm}[Stavrova, 2018]
  Let $R$ be a Dedekind domain of arithmetic type (e.g. $R=\mathbb{Z}$), and let $G$ be a simply connected Chevalley--Demazure group scheme of rank $\ge 2$. Then \( K_1^G(R[x_1,\ldots,x_n]) = \{ * \}. \)
 \end{thm}
\end{frame}

\begin{frame}{Unstable $K_1$-functor of an isotropic reductive group}

Let $k$ be a field and let $G$ be an isotropic reductive group (so contains a proper parabolic $k$-subgroup $P$).
\pause

Let $P \subset G$ be a proper parabolic subgroup and let $P^{-}$ be an opposite
parabolic subgroup. Denote by \( U_P,\) \(U_{P^-} \) their unipotent radicals.
For any $k$-algebra $A$ define
\[ E_P(A)=\langle U_P(A),\,U_{P^-}(A)\rangle \subset G(A),\ K_1^{G,P}(A) := G(A)/E_P(A).\]
\pause

If $G$ has isotropic rank $\ge 2$, then $E_P(A)$ does not depend on $P$.
\pause
\begin{exm}
 For $G=\mathrm{GL}_n$ split, \( E_P(A) = E_n(A) \) and \( K_1^G(A)=\mathrm{GL}_n(A)/E_n(A). \)
\pause
$K_1^{GL_n}(A)$ is trivial iff $\mathrm{GL}_n(A) = E_n(A)$.
\end{exm}
\end{frame}

\section{Bass--Quillen for $\K_2$}

\begin{frame} \tableofcontents[currentsection] \end{frame}

\begin{frame}{Analogue of Bass-Quillen for $\K_2$?}
Compute $\mathcal{H}_{k,*}^{\mathbb{A}^1}(S^1_s \wedge X_+, G) = ?$
\pause

\begin{thm}[Tulenbaev 1983]
  For an arbitrary Noetherian $A$ the canonical stabilization homomorphism
  $\K_2^{SL_r}(A[X_1, \ldots X_n]) \to \K_2(A[X_1, \ldots X_n])$
  \begin{itemize}
   \item is surjective, provided $r \geq \mathrm{max}(4, \mathrm{dim}(A) + 2)$;
   \pause
   \item is injective, provided $r \geq \mathrm{max}(5, \mathrm{dim}(A) + 3)$.
  \end{itemize}
\end{thm}
\pause


\begin{thm}[V\"olkel--Wendt 2016]
 Assume $k$ - infinite, then for $G$ of isotropic rank $\geq 2$ one has $\mathcal{H}_{k,*}^{\mathbb{A}^1}(S^1_s \wedge \mathrm{Spec}(k)_+, G) = \mathrm{H}_2(G(k), \mathbb{Z})$. (?)
\end{thm}

\end{frame}

\begin{frame}{Motivic Matsumoto}
 \begin{thm} [Morel--Sawant 2022]
If $G = G_{sc}(\Phi, -)$ is a Chevalley--Demazure groups scheme, then there exists an isomorphism of sheaves
\[ \pi_{1}^{\mathbb{A}^1}(G) \cong
\begin{cases}
\K_2^{M}, & \text{if $G$ is not of symplectic type};\\
\K_2^{MW}, & \text{if $G$ is of symplectic type}.
\end{cases} \]
\pause
\[ \pi_1^{\mathbb{A}^1}(G) \;:=\; a_{\mathrm{Nis}} \Big( X \;\longmapsto\; \mathcal{H}^{\mathbb{A}^1}_ {k, *}(S^1_s \wedge X_+,\, G) \Big), \]
\pause
\[ \mathbf{K}_2^{M(W)} \;:=\; a_{\mathrm{Nis}} \Big( U \;\longmapsto\; K_2^{M(W)}(\mathcal{O}(U)) \Big). \]
\end{thm}
\end{frame}

\begin{frame}{Milnor and Milnor--Witt $K_2$ Sheaves}

Set $A = \mathcal{O}(U)$. One has \[ K_2^M(A) = \frac{A^\times \otimes_{\mathbb Z} A^\times}
{\langle a \otimes (1-a) \mid a,1-a \in A^\times \rangle}. \]
\pause

The graded Milnor--Witt $K$-ring $K_*^{MW}(A)$ is the $\mathbb Z$-graded ring
generated by symbols $[u]$ for $u\in A^\times$ in degree $1$ and an element
$\eta$ in degree $-1$, subject to the relations
\[
[uv]=[u]+[v]+\eta[u][v],\
[u][1-u]=0,\
\eta[u]=[u]\eta,\
\eta(2+\eta[-1])=0.
\]
\pause
Then
\( K_2^{MW}(A) := K_*^{MW}(A)_2. \)
\pause

\[ \mathbf{K}_2^{MW}(\mathcal{O}(U)) = \bigcap_{x\in U^{(1)}} \ker\!\left( \partial_x : K_2^{MW}(k(U)) \to K_1^{MW}(\kappa(x)) \right), \]
i.e.\ the subgroup of \emph{unramified elements} in $K_2^{MW}(k(U))$.
\end{frame}

\begin{frame}{Analogue of Bass-Quillen for $\K_2$}
Now let $G$ be isotropic alg. gp over $k$. For $\alpha \in \Phi_P$ there exists \( X_\alpha : V_\alpha \to G, \) where $V_\alpha$ is a f.g. proj. $k$-module.
\pause
By definition, the Steinberg group $\St^G(R)$ is presented by  relative root elements $X_\alpha(v)$, $v \in V_\alpha$ and relations

\( [X_\alpha(u),X_\beta(v)] = \prod_{i,j>0} X_{i\alpha+j\beta}\big(N_{\alpha\beta ij}(u,v)\big), \)
where $N_{\alpha\beta ij}$ are polynomial maps \( V_\alpha \times V_\beta \to V_{i\alpha+j\beta}, \alpha,\beta \in \Phi_P.\)
\pause

Set $\K_2^G(R) := \mathrm{Ker}(\St^G(R) \to E(R) \subseteq G(R))$.
\pause

 \begin{conj}[Bass--Quillen for $\K_2$] $G$ is of isotropic rank $\geq 3$, and $R$ regular, then $\K_2^G(R[x]) = \K_2^G(R)$ and
 $\mathcal{H}_{k,*}^{\mathbb{A}^1}(S^1_s \wedge (\mathrm{Spec}(R))_+, G) = \K_2^G(R[x])$. \end{conj}
\end{frame}

\begin{frame}{Analogue of Bass--Quillen for $\K_2$}
\begin{thm}[Lavrenov-S.-Voronetsky 2022 + S. 2024]
 $G$ -- Chevalley--Demazure of type $\Phi = \mathsf{A}_{\geq 4}, \mathsf{D}_{\geq 5}, \mathsf{E}_6, \mathsf{E}_7$, $R$ regular algebra over a field $k$. Then BQ for $\K_2$ holds.
\end{thm}
\pause

\begin{thm}[S. 2024] Under the same assumption on $G$ but for $R$ - a Dedekind domain one has $\K_2(R[x_1, \ldots, x_n]) = \K_2(R)$.
\end{thm}
\end{frame}

\section{Proof: first reduction}

\begin{frame} \tableofcontents[currentsection] \end{frame}

\begin{frame}{A sufficient condition for $\mathbb{A}^1-$invariance}
Let $k$ be a field and let $K\colon \catname{Alg}_k \to \catname{Groups}$ be a functor satisfying:
\begin{enumerate}
  \item $K$ commutes with filtered colimits;
  \pause
  \item For $R \in \catname{Alg}_k$ and $a\in R$ consider the pull-back square
  \[ \begin{tikzcd}[ampersand replacement=\&, column sep=1.1em]
                                                 \&                      \&\& \only<5->{\Ker(K(e))\ar[hookrightarrow]{r} \ar[hookrightarrow]{d}} \& \only<5->{\Ker(K(ev(0))) \ar[hookrightarrow]{d} } \\
    R \ltimes t R_a[t] \ar{r}{l} \ar{d}[swap]{e} \ar[phantom, very near start]{dr}{\lrcorner} \& R_a[t] \ar{d}{ev(0)} \&\& \only<5->{K(R\ltimes t R_a[t]) \ar{r} \ar{d} \& K(R_a[t]) \ar{d}}  \\
    R \ar{r}{\lambda_a}                          \& R_a.                 \&\& \only<5->{K(R)\ar{r}                  \& K(R_a)}
  \end{tikzcd}\]
  Then the application of $K$ induces an injection on kernels
 \end{enumerate}
\end{frame}

\begin{frame}{A sufficient condition for $\mathbb{A}^1-$invariance}
 \begin{enumerate}
 \pause \setcounter{enumi}{2}
  \item $K$ (weak Nisnevich excision) Let $\iota \colon B \hookrightarrow A$ be an injective homomorphism of domains, contained in
  $\catname{Alg}_k$ such that $\mathrm{Spec}(A)\to\mathrm{Spec}(B)$ is an etale morphism of affine schemes. Let $h \in B$ be such that $\iota$ induces an isomorphism $B / hB \cong A / \iota(h)A$. \pause Consider the dual Nisnevich square:
  \scalebox{0.9}{\begin{minipage}{1.00\linewidth}\[\begin{tikzcd}[ampersand replacement=\&, column sep=1.2em]
                           \&          \& \only<4->{\Ker(K(\lambda_h)) \ar{d} \ar[twoheadrightarrow]{r} \& \Ker(K(\lambda_{\iota(h)})) \ar{d}}  \\
    B \ar{r}{\iota} \ar{d}{\lambda_h} \ar[phantom, very near start]{dr}{\lrcorner} \& A \ar{d}{\lambda_{\iota(h)}} \& \only<4->{K(B)  \ar{r}\ar{d} \& K(A) \ar{d}}    \\
    B_h \ar{r}             \& A_{\iota(h)}      \& \only<4->{K(B_h)\ar{r} \& K(A_{\iota(h)})} \end{tikzcd}\]\end{minipage}}

  Then the application of $K$ induces a surjection on kernels.
 \end{enumerate}
\end{frame}

\begin{frame}{A sufficient condition for $\mathbb{A}^1-$invariance}
 \begin{enumerate}
 \pause \setcounter{enumi}{3}
  \item ($\mathbb{P}^1$-glueing, Horrocks theorem) For any local domain $R \in \catname{Alg}_k$ the following diagram is a pull-back square:
  \[\begin{tikzcd}[ampersand replacement=\&, column sep=1.2em]
  K(R) \ar[phantom, very near start]{dr}{\lrcorner} \ar[r] \ar[d] \& K(R[t]) \arrow{d} \\ K(R[t^{-1}]) \ar{r} \& K(R[t, t^{-1}]). \end{tikzcd}\]
  \pause \item For any field extension $E/k$ one has $K(E[t]) = K(E).$
 \end{enumerate}
 \pause Then for any geom. regular $R$ over $k$ one has $K(R[t]) = K(R)$.
\end{frame}

\begin{frame}{Horrocks theorem}
\begin{thm}[Horrocks 1964]  Let $R$ be a commutative local ring. If a vector bundle $B$ on $\mathbb{A}^1_R$ extends to a vector bundle on $\mathbb{P}^1_R$, then $B$ is a trivial bundle. \end{thm}
\end{frame}

\begin{frame}{Horrocks theorem for $K_1$}
Let $A$ be a comm. local ring with residue field $k = A/M$.
\pause

 Define the \textit{relative elementary group}:
 $\E(n, A[X^{\pm 1}], M[X^{\pm 1}]) = \mathrm{Ker}(\E(n, A[X^{\pm 1}]) \to \E(n, k[X^{\pm 1}])).$
\pause

 Key fact in the proof $K_1$-analogue:
\begin{prop}[Suslin 1977]
 \( \E(n, A[X^{\pm 1}]) = \E(n, A[X]) \cdot B(A[X^{\pm 1}]) \cdot \E(n, A[X^{\pm 1}], M[X^{\pm 1}]). \)
\end{prop}
\pause

 \textbf{Sketch:} Construct a family of conjugation automorphisms $\delta_i \in \mathrm{Aut}(E(n, A[X^{\pm 1}]))$.
\pause
 Each $\delta_i$ models conjugation by $diag(1, \ldots 1, X, 1, \ldots, 1)$ automorphism.
\pause
 These commute and stabilize the decomposition in the right-hand side.

\end{frame}

\begin{frame}{Horrocks theorem for $K_2$}
 Let $(A, M)$ be a comm. local ring with residue field $k$.
\pause

 Set $B := A[X^{-1}] + M[X]$, $R := A[X, X^{-1}].$
\pause

 Prove that both homomorphisms are injective:
 \[ \St(\Phi, A[X^{-1}]) \xrightarrow{j_B} \St(\Phi, B) \xrightarrow{j_R} \St(\Phi, R), \]
\pause

 \begin{thm}[Lavrenov--S. 2019] Let $\Phi$ be a simply-laced root system containing $\mathsf{A_4}$.
 Assume that $j_R$ is injective. Then $j_B$ is injective as well.
 \end{thm}
\pause

  \textbf{Sketch:} Set $G^{\geq 0}_M := \mathrm{Im}\left(\St(\Phi, A[X], M[X]) \to \St(\Phi, R)\right). $
\pause
  $\St(\Phi, B)$ acts simply transitively on the set of triples
  \[G^{\geq 0}_M \times \St(\Phi, A[X^{-1}]) \times (1 + M)^\times.\]
\end{frame}

\begin{frame}{Horrocks theorem for $K_2$}
The injectivity of $j_R$ is known in the following cases:
 \begin{itemize}
  \item $\Phi=\mathsf{A}_{\geq 4}$. Tulenbaev 1983, direct proof;
  \pause
  \item $\Phi=\mathsf{D}_{\geq 7}$, require additionally that $2 \in A^\times$. Lavrenov-S. 2019, uses Panin's stabilization for orthogonal K-theory combined with Bass Fundamental Theorem for Grothendieck--Witt groups proved by J. Hornbostel in 2005.
  \pause
  \item $\Phi = \mathsf{A}_{\geq 4}, \mathsf{D}_{\geq 5}, \mathsf{E}_6, \mathsf{E}_7$ under the additional assumption that $A$ is a \textbf{domain}, S. 2024. \textbf{Idea}: Tulenbaev's 1983 direct proof can be generalized.
 \end{itemize}
\end{frame}

\begin{frame}{Relative Steinberg groups}
 Let $R$, $I \trianglelefteq R$ be arbitrary. By their definition, relative Steinberg groups are certain central extensions:
 \[ C(\Phi, R, I) \hookrightarrow \St(\Phi, R, I) \xrightarrow{\mu} \St(\Phi, R) \rightarrow \St(\Phi, R/I). \]
\pause

 If $I$ is a splitting ideal of $R$ (e.g. $\langle X \rangle \trianglelefteq A[X]$), then $C(\Phi, R, I) = 1$.
\pause

 Now let $(A, M)$ be local.
\pause

 $B = A[X^{-1}] + M[X]$, $R = A[X^{\pm 1}]$, $I = M[X^{\pm 1}]$.
\pause

 Simple diagram chasing reduces the injectivity of $j_R$ to the surjectivity of the right arrow:
 \[C(\Phi, A[X], M[X]) \longrightarrow C(\Phi, B, I) \longrightarrow C(\Phi, R, I).\]
\pause

 We will prove that the composite arrow is surjective.
\end{frame}

\begin{frame}{Analogue of Suslin's decomposition for $\K_2$}
  Need an analogue of Suslin's decomposition. Simply proving an analogue of this decomposition on the level of Steinberg groups is not enough:
\pause

  \[ \St(n, A[X^{\pm 1}]) = \St(n, A[X]) \cdot \widetilde{B}(A[X^{\pm 1}]) \cdot \St(n, A[X^{\pm 1}], M[X^{\pm 1}]). \]
\pause

  Need something much more refined, something that e.g. takes care of the uniqueness of decompositions in Suslin's structure theorem.
\end{frame}

\section{Proof: a more deep dive}

\begin{frame} \tableofcontents[currentsection] \end{frame}

\begin{frame}{Crossed modules}
\begin{itemize}
\item Let $N, M$ be groups, $N \curvearrowright N$ by right conjugation, $N \curvearrowright M$ on the right.
A $\mu\colon M \to N$ is a \textit{precrossed module} if $\mu$ preserves the action of $N$, i.\,e.
\( \mu(m^n) = \mu(m)^n, \text{for all $m \in M$, $n\in N;$} \)
\pause

\item Let $\mu\colon M \to N$ and $\mu' \colon M' \to N'$ be a pair of precrossed modules.
A \textit{homomorphism of precrossed modules} $(f, g)\colon \mu \to \mu'$ is a pair of group homomorphisms
$f\colon M \to M'$, $g\colon N \to N'$ such that $\mu'f = g \mu$ and ${f(m)}^{g(n)} = f(m^n)$ for $n \in N$, $m \in M$.
\pause

\item $\mu$ is called a \textit{crossed module} if \({m}^{\mu(m')} = {m'}^{-1} m m'\), for all $m, m' \in M.$
\pause

\item Base example: $\mu \colon \St(\Phi, R, I) \to \St(\Phi, R)$.
\end{itemize}
\end{frame}


\begin{frame}[fragile]{Formalism of triples}

\begin{columns}[T]

\begin{column}{0.55\textwidth}
\centering
\[
\begin{tikzcd}[row sep=2.6em, column sep=3em]
G_{123}
  \arrow[rr, "f_{23}"]
  \arrow[dd, "f_{12}"']
  \arrow[rd, "f_{13}"]
&&
G_{23}
  \arrow[dd, dashed, "{g_2^3}" description, pos=0.3]
  \arrow[rd, "{g_3^2}", pos=0.3]
& \\
&
G_{13}
  \arrow[rr, "{g_3^1}" description, pos=0.3]
  \arrow[dd, "{g_1^3}" description, pos=0.3]
&&
G_3
  \arrow[dd, "h_3"]
\\
G_{12}
  \arrow[rr, dashed, "{g_2^1}" description, pos=0.3]
  \arrow[rd, "{g_1^2}", pos=0.3]
&&
G_2
  \arrow[rd, dashed, "h_2"']
& \\
&
G_1
  \arrow[rr, "h_1"']
&&
G
\end{tikzcd}
\]
\end{column}

\begin{column}{0.45\textwidth}

\begin{itemize}
\item $G_1 \curvearrowright G_{13}$ on the right.
\pause
\item $G \curvearrowright G_3$ on the right.
\pause
\item $g_1^3$ is a precrossed module.
\pause
\item $h_3$ is a crossed module.
\pause
\item $(g_3^1,h_1)\colon g_1^3 \to h_3$ is a morphism of precrossed modules.
\end{itemize}

\end{column}

\end{columns}
\end{frame}

\begin{frame}{Formalism of triples, continued}
Set $V = G_1 \times G_2 \times G_3$, $W = G_{12} \times G_{13} \times G_{23}$.
\pause

Define star operation $\star \colon W \times V \to V$ as follows.
For $v = (x, y, z) \in V$ and $w = (a, b, c) \in W$ set
\begin{multline*} (a, b, c) \star (x, y, z) = (x \cdot g_1^3(b) \cdot g_1^2(a),\\ g_2^1(a)^{-1} \cdot y \cdot g_2^3(c),\ g_3^2(c)^{-1} \cdot g_3^1(b)^{-h_2(y)} \cdot z).\end{multline*}
\pause

\begin{equation*}(a', b', c') \star \left( (a, b, c) \star v \right) = (a \cdot a', b \cdot {b'}^{g_1^2(a)^{-1}}, c \cdot c') \star v.\end{equation*}
\pause

Let $h \colon V \to G$ be given by $(x, y, z) \mapsto h_1(x) \cdot h_2(y) \cdot h_3(z)$, $h$ gives rise to a well-defined map $\overline{h} \colon \overline{V} \to G$,
where $\overline{V}$ denotes the set of equivalence classes $V/\sim_W.$ We use the notation $[a, b, c]$ to denote the equivalence class of the triple $(a, b, c) \in V$.


\end{frame}

\begin{frame}{Formalism of triples, continued}
\begin{lemma} Assume that the back face $(f_{12}, f_{23}, g_2^1, g_2^3)$ in the cube diagram is a pullback square.
Let $z \in G_3$ be such that $[1, 1, 1] = [1, 1, z]$.
Then $z \in g_3^1(\Ker(g_1^3)).$ \end{lemma}
\pause
\textbf{Proof:} There exists $(a, b, c)\in W$ such that
\begin{multline*} (a, b, c) \star (1, 1, 1) = ( g_1^3(b) \cdot g_1^2(a),\\ g_2^1(a)^{-1} \cdot g_2^3(c),\ g_3^2(c)^{-1} \cdot g_3^1(b)^{-1}) = (1,1,z). \end{multline*}
\pause
By assumption there exists $e \in G_{123}$ such that $f_{12}(e) = a$, $f_{23}(e) = c$, hence
\begin{multline*} 1 = g_1^3(b) \cdot g_1^2(f_{12}(e)) = g_1^3(b \cdot f_{13}(e)),\\ z = g_3^2(f_{23}(e))^{-1} \cdot g_3^1(b)^{-1} = g_3^1(b \cdot f_{13}(e))^{-1}. \end{multline*}
\end{frame}

\begin{frame}[fragile]{The cube}

\begin{columns}[T]
\begin{column}{0.4\textwidth}
\raggedleft
\scriptsize
\hspace*{0pt}%
\[
\begin{tikzcd}[row sep=1.8em, column sep=2.1em]
U^+_M
  \arrow[rr, hook]
  \arrow[dd, hook]
  \arrow[rd, hook]
&&
B_M
  \arrow[dd, dashed]
  \arrow[rd, hook]
& \\
&
G_M^+
  \arrow[rr]
  \arrow[dd, "\mu", pos=0.3]
&&
G_M
  \arrow[dd, "\mu"]
\\
U^+
  \arrow[rr, hook, dashed]
  \arrow[rd, hook]
&&
B
  \arrow[rd, hook, dashed]
& \\
&
G^+
  \arrow[rr]
&&
G
\end{tikzcd}
\]
\end{column}

\begin{column}{0.6\textwidth}
\raggedleft
\scriptsize
\[
\begin{aligned}
G \;&=\; \St(\Phi, A[X,X^{-1}]),\\
G^+ \;&=\; \St(\Phi, A[X]),\\
B \;&=\; \UU(\Phi^+, A[X,X^{-1}])
       \rtimes \StH(\Phi, A[X,X^{-1}]) \le G,\\
G_M \;&=\; \St(\Phi, A[X,X^{-1}], M[X,X^{-1}]),\\
U^+ \;&=\; \UU(\Phi^+, A[X]),\\
G_M^+ \;&=\; \St(\Phi, A[X], M[X]),\\
B_M \;&=\; \UU(\Phi^+, M[X,X^{-1}]) \times \{X,1+M\} \le G_M,\\
U_M^+ \;&=\; \UU(\Phi^+, M[X]).
\end{aligned}
\]
\end{column}
\end{columns}
\medbreak

\scriptsize { Here $\StH(\Phi, A[X^{\pm 1}]) = \langle h_\alpha(u) \mid u \in A[X^{\pm 1}]^\times, \alpha \in \Phi \rangle \leq \St(\Phi, A[X, X^{-1}])$. }
\pause

\scriptsize { The arrow $B_M \to B$ is injective since $(1 + M)^\times \to G$ is injective for \textbf{domains}. }
\end{frame}

\begin{frame}{The main result}
\begin{thm}
 The group $G = \St(\Phi, A[X, X^{-1}])$ acts on $V/\sim_W$, this action is compatible with the action of $G^+$ and, moreover, for every $z \in G_M$ one has $\mu(z) \cdot [1, 1, 1] = [1, 1, z]$.
\end{thm}
\pause

\begin{cor}
 $C(\Phi, A[X], M[X]) \twoheadrightarrow C(\Phi, A[X, X^{-1}], M[X, X^{-1}])$.
\end{cor}
\pause
\textbf{Proof:} Take $z$ from RHS. Then $\mu(z) = 1$, hence by the above lemma $z$ comes from $\Ker(G_M^+ \to G^+)$ via $G_M^+ \to G_M$.
\end{frame}

\begin{frame}{Weight automorphisms}
 If we are to repeat the idea of Tulenbaev we want to construct analogues of ``weight'' automorphisms $\delta_\varpi$ on the level of $V/\sim_W$.
\pause
Constructing the action of $G$ is hard, we will construct the action of some extension of $G$, for which it is easier.

 \medbreak

 To $\omega \in P(\Phi)$ and $\beta \in \mathbb{Z} \Phi$ the product $\langle \omega, \beta\rangle \in \mathbb{Z}$.
\pause

 For $u \in R^\times$ and $\omega \in P(\Phi)$ there exists a bijection of the generating set $\St(\Phi, R)$ via the following mapping:
 \[ x_\alpha(a) \mapsto x_\alpha(u^{\langle\omega, \alpha\rangle} \cdot a),\ \alpha\in \Phi,\ a \in R. \]
\pause

 This gives rise to an automorphism $\chi_{\omega, u}$ of $\St(\Phi, R)$.
\end{frame}

\begin{frame}{Certain presentation $\St(\Phi, A[X, X^{-1}])$.}
Let $\Phi= \mathsf{D}_\ell, \mathsf{E}_6$ or $\mathsf{E}_7$.
Let $k$ be the green vertex on the diagram:

\medbreak
\begin{tabular}{ccc}
\scalebox{0.52}{%
\begin{tikzpicture}
\node[root, highlighted, label=$1$] (d1) at (0,0.5) {};
\node[root, label=$2$] (d2) at (1,0) {};
\node[root, label=$3$] (d3) at (2,0) {};
\node at (3,0) {\ldots};
\node[root, label={[xshift=-0.5cm,yshift=-0.6cm]$\ell-2$}] (dl2) at (4,0) {};
\node[root, label={[xshift=-0.6cm,yshift=-0.35cm]$\ell-1$}] (dl1) at (5,0.5) {};
\node[root, label={[xshift=-0.5cm,yshift=-0.3cm]$\ell$}] (dl) at (5,-0.5) {};
\node[root, zeroroot, label=below:{$0$}] (d0) at (0,-0.5) {};

\draw (d1) -- (d2) -- (d3) -- (2.5,0);
\draw (3.5,0) -- (dl2) -- (dl1);
\draw (dl2) -- (dl);
\draw[dottededge] (d2) -- (d0);

\node[levi, fit=(d2) (d3) (dl) (dl1) (dl2), label=above:{$\Delta$}] {};
\end{tikzpicture}}
&
\scalebox{0.52}{%
\begin{tikzpicture}
\node[root, highlighted, label=$1$] (e1) at (0,0) {};
\node[root, label=$3$] (e3) at (1,0) {};
\node[root, label={[xshift=-0.3cm]$4$}] (e4) at (2,0) {};
\node[root, label=$5$] (e5) at (3,0) {};
\node[root, label=$6$] (e6) at (4,0) {};
\node[root, label={[xshift=-0.3cm,yshift=-0.3cm]$2$}] (e2) at (2,1) {};
\node[root, zeroroot, label=left:$0$] (e0) at (2,2) {};

\draw (e1) -- (e3) -- (e4) -- (e5) -- (e6);
\draw (e2) -- (e4);
\draw[dottededge] (e2) -- (e0);

\node[levi, fit=(e2) (e3) (e4) (e5) (e6), label=above:{$\Delta$}] {};
\end{tikzpicture}}
&
\scalebox{0.52}{%
\begin{tikzpicture}
\node[root, zeroroot, label={$0$}] (f0) at (0,0) {};
\node[root, label={$1$}] (f1) at (1,0) {};
\node[root, label={$3$}] (f3) at (2,0) {};
\node[root, label={[xshift=-0.3cm]$4$}] (f4) at (3,0) {};
\node[root, label={$5$}] (f5) at (4,0) {};
\node[root, label={$6$}] (f6) at (5,0) {};
\node[root, label={[xshift=-0.3cm,yshift=-0.3cm]$2$}] (f2) at (3,1) {};
\node[root, highlighted, label={$7$}] (f7) at (6,0) {};

\draw (f0) -- (f1) -- (f3) -- (f4) -- (f5) -- (f6) -- (f7);
\draw (f2) -- (f4);

\node[levi, fit=(f1) (f2) (f3) (f4) (f5) (f6), label=above:{$\Delta$}] {};
\end{tikzpicture}}
\\[-0.2em]
$\mathsf{D}_\ell$ & $\mathsf{E}_6$ & $\mathsf{E}_7$
\end{tabular}
\medbreak
\pause

$m_k(\alpha_\mathrm{\max}) = 1$, therefore $m_k(\alpha) = 1$ for all $\alpha \in \Sigma_k^+$ and the subgroup $\UU(\Sigma^+_k, R)$ is abelian.
\pause
Also the weight $\varpi_k$ is a \textit{microweight} of $\Phi$.
\end{frame}


\begin{frame}{Certain presentation $\St(\Phi, A[X, X^{-1}])$.}
Denote by $G'$ the amalgamated product of $\St(\Phi, A[X])$ and the cyclic group $\langle \sigma \rangle$ amalgamated over the subgroup generated by the following relations:
\begin{align}
{}^\sigma x_{\alpha}(f) = & x_{\alpha} (Xf), & \alpha \in \Sigma^+_k, f \in A[X], \label{eq:sigma-sigma-plus} \\
x_{\beta}(f)^ \sigma     =& x_{\beta} (Xf), & \beta \in \Sigma^-_k, f \in A[X], \label{eq:sigma-sigma-minus} \\
[\sigma,\, x_\gamma(f)]   =& 1, & \gamma \in \Delta_k, f \in A[X]. \label{eq:sigma-delta}
\end{align}

\end{frame}

\begin{frame}{Certain presentation $\St(\Phi, A[X, X^{-1}])$.}

\begin{block}{Proposition}
For $\Phi$ and $G'$ as above there exists a group homomorphism
\[
\varphi \colon \St(\Phi,A[X,X^{-1}]) \to G'
\]
making the following diagram commute:
\[
\begin{tikzcd}[ampersand replacement=\&, column sep=2.8em, row sep=2.2em]
\& \St(\Phi,A[X,X^{-1}]) \arrow[rd, dashed, "\varphi"] \& \\
\St(\Phi,A[X]) \arrow[ru] \arrow[rr] \&\& G'
\end{tikzcd}
\]
\end{block}
\end{frame}

\begin{frame}{Construction of $\delta_\varpi$}

Consider the subset $\mathcal{X}_\omega$ of the generating set of $\St(\Phi, A[X])$ consisting of those generators $x_{\alpha}(a) \in \St(\Phi, A[X])$ for which
$\langle \omega, \alpha\rangle < 0$ implies that $a \in A[X]$ is divisible by $X^{-\langle \omega, \alpha\rangle}$.
\pause

If $\langle \omega, \alpha\rangle \geq 0$ then $\mathcal{X}_\omega$ contains $x_\alpha(a)$ for all $a \in A[X]$.
\pause

Denote by $N_{\omega}$ the subgroup of $\St(\Phi, A[X])$ generated by $\mathcal{X}_\omega$.
\pause

Let $\omega \in P(\Phi)$ be a weight. By definition, an \textit{$\omega$-pair} is a pair of mutually inverse group homomorphisms
\( N_\omega \overset{\delta_\omega}{\underset{\delta_{-\omega}}{\rightleftarrows}} N_{-\omega} \) satisfying
\[ \delta_{\pm\omega}(x_\alpha(f)) = x_\alpha\!\bigl(X^{\pm\langle\omega,\alpha\rangle}\cdot f\bigr), \qquad x_\alpha(f)\in \mathcal{X}_{\pm\omega}. \]
\end{frame}

\begin{frame} [fragile]{The $\sigma$-diagram}

\[
\begin{tikzcd}[ampersand replacement=\&, column sep=3.2em, row sep=2.6em]
N_\omega
  \arrow[r, "\delta_\omega"]
  \arrow[d, hook]
  \arrow[rr, bend left=30, "\mathrm{id}"]
\&
N_{-\omega}
  \arrow[r, "\delta_{-\omega}"]
  \arrow[d, hook]
\&
N_\omega
  \arrow[d, hook]
\\
\St(\Phi,A[X])
  \arrow[d]
\&
\St(\Phi,A[X])
  \arrow[d]
\&
\St(\Phi,A[X])
  \arrow[d]
\\
\St(\Phi,A[X^{\pm1}])
  \arrow[r, "\chi_{\omega,X}"']
  \arrow[rr, bend right=30, "\mathrm{id}"']
\&
\St(\Phi,A[X^{\pm1}])
  \arrow[r, "\chi_{-\omega,X}"']
\&
\St(\Phi,A[X^{\pm1}])
\end{tikzcd}
\]
\end{frame}

\begin{frame}
 \begin{dfn} \label{sigma-def}
Let $k$ be the green vertex, let $\omega = \pm \varpi_k$.
Denote by $V_\omega$ the subset of $V$ consisting of those triples $v = (x, y, z)$ for which $x \in N_\omega\cdot \mathrm{W}(\Phi, A)$.
We define the function $\delta_\omega \colon V_\omega \to V_{-\omega}$ via the following formula:
\begin{equation} \label{eq:sigma-def} (x_1 \cdot x_2, y, z) \mapsto (\delta_\omega(x_1)\cdot x_2, x_2^{-1} \cdot \chi_{\omega, X}(x_2 \cdot y),\  \widetilde{\chi}_{\omega, X}(z)), \end{equation}
where $x_1 \in N_\omega$, $x_2 \in \mathrm{W}(\Phi, A)$,\ $y \in B$, $z \in G_M$, and
$\chi_{\omega, X}$, $\widetilde{\chi}_{\omega, X}$ denote the ``weight'' automorphism and its relative analogue.
\end{dfn}
\end{frame}

\begin{frame}
 It is possible to prove the correctness of this definition using the following:
\pause

 \begin{prop} \label{prop:v-correctness2}
For $\omega = \pm \varpi_k$ and every $w \in W$, $v \in V_\omega$ such that $v' = w \star v \in V_\omega$ there exists $w' \in W$ such that
$\delta_\omega(w \star v) = w' \star \delta_\omega(v)$.
\end{prop}
\pause

There is well-defined $\overline{\delta}_{\pm \omega} \colon V/\sim_W \to V/\sim_W$.
\pause

$G'$ acts on $V/\sim_W$: $\sigma$ acts via $\overline{\delta}_{\pm \omega}$, $\St(\Phi, A[X])$ acts on the first component of $V$ via left multiplication.
\pause

The definining relations of $G'$ are verified. Consequently $G'$ acts on $V/\sim_W$.

\end{frame}


\begin{frame}
\[ \text{Thank you for attention!} \]
\end{frame}

\end{document}
